% ==============================================================================
% 5. CONSIDERAÇÕES FINAIS
% ==============================================================================
\chapter{Considerações finais}
\label{cap:consideracoes}

Este trabalho perguntou quão bons são os testes que uma LLM gera quando se pede a ela, numa única mensagem, que teste um módulo. Em vez de responder com cobertura ou com uma comparação contra testes humanos cuja qualidade não é medida, respondemos com a distância até um teto: a melhor suíte que conseguimos construir para o alvo, com um desenvolvedor, LLMs de fronteira e ferramentas de mutação e de \textit{smells}, e na qual cada mutante que sobrevive tem prova de que não pode ser morto. Para um avaliador de expressões com 369 mutantes, esse teto ficou em 99,5\% de escore de mutação, com cobertura de 100\% de linhas e 99,3\% de ramos, nenhum \textit{smell} e dois mutantes com prova de equivalência.

Contra esse teto, as respostas são estas. \textbf{QP1:} com o prompt básico, as suítes geradas ficam a 8,4 e 8,7 pontos percentuais do teto em escore de mutação no caso dos dois modelos de fronteira e a 17,9 no caso do modelo de baixo custo, enquanto em cobertura ficam a 3 pontos ou menos em linhas, em precisão a 2 pontos ou menos e em \textit{smells} a menos de 0,1 por teste. \textbf{QP3:} a dimensão em que a distância é maior é sempre a de detecção de defeitos, por margem de duas a sete vezes sobre a cobertura. \textbf{QP2:} informar os critérios de avaliação no prompt reduziu a distância em detecção em 57\%, 75\% e 73\%, em todas as gerações de todos os modelos, a ponto de o modelo de baixo custo orientado superar os dois modelos de fronteira sem orientação; a melhor condição, Gemini com P2, chegou a 2,2 pontos do teto; o custo foram suítes maiores e alguns falsos positivos por asserções exatas demais, dois dos quais apontavam um defeito real do alvo. \textbf{QP4:} a dispersão entre gerações depende do modelo e cai com a orientação: meio ponto nos modelos de fronteira orientados, para os quais uma geração basta, e 2,5 pontos no modelo de baixo custo com o prompt básico, para o qual seriam necessárias sete; uma geração ao acaso por condição acerta sempre as conclusões grandes (prompt, grupo de modelos) e erra entre 4\% e 36\% das vezes a ordem entre condições vizinhas.

Dois resultados valem independentemente dos números das suítes geradas. O primeiro é metodológico: pontuar mutação em paralelo inflou o escore do próprio teto em 2,4 pontos percentuais por \textit{timeouts} que não eram detecções, e essa inflação não é uniforme entre suítes, o que compromete comparações; toda pontuação com limite de tempo adaptativo deveria ser sequencial, e mutantes mortos e por \textit{timeout} deveriam ser reportados em separado. O segundo é sobre o próprio conceito de teto: uma suíte que parecia saturada ainda deixava 19 mutantes matáveis em fronteiras de classes de caracteres e em mensagens de erro, e a análise que se seguiu ainda classificou como inalcançáveis seis mutantes que uma suíte gerada depois alcançou por injeção de dependência. O que transformou a suíte em referência foi a análise individual dos sobreviventes e a revisão feita quando uma suíte gerada a superou. Uma referência de qualidade precisa de prova e de um critério explícito para ser revisada.

As limitações são as de um estudo com um alvo, três modelos e cinco gerações por condição: os números são específicos do alvo e do momento, e a dispersão foi medida com amostras pequenas. Os próximos passos seguem daí. Primeiro, repetir o protocolo em outros alvos, com características diferentes (E/S, estado, dependências), para ver se a distância e sua estrutura se mantêm. Segundo, acrescentar uma condição em que o modelo pode executar os testes e receber o resultado, para medir quanto do caminho até o teto o laço de execução fecha sozinho; se fechar quase tudo, a conclusão prática muda. Terceiro, avaliar a seleção entre várias gerações (gerar $k$ vezes e ficar com a melhor pelo escore de mutação) como estratégia de uso, com seu custo. Por fim, o aparato e o teto estão públicos e podem servir de \textit{benchmark} para modelos futuros: a distância até o teto é uma medida que não depende de quem escreveu a suíte humana de comparação, e por isso pode ser acompanhada ao longo do tempo.
